\documentclass[11pt]{article}
\usepackage[margin=1in]{geometry}
\usepackage{booktabs}
\usepackage{graphicx}
\usepackage{hyperref}
\usepackage{xcolor}
\hypersetup{colorlinks=true,linkcolor=black,citecolor=black,urlcolor=blue}

\title{Access-Controlled Text-to-SQL Without Showing the Model the Restricted Schema:\\
Results on Two Public Benchmarks, and a Fail-Closed Bug They Exposed}
\author{Ashish Sinha\\ \small Independent \\ \small \url{https://github.com/ashishsinha1602/schemagate}}
\date{October 2026}

\emergencystretch=3em
\begin{document}
\maketitle

\begin{abstract}
Text-to-SQL systems deployed inside organisations must answer the questions a user may ask and refuse
the ones they may not. Two recent benchmarks measure exactly this: the RBAC Text-to-SQL benchmark of
Fei et al.\ and Spider-ACL of Klisura et al., whose labels come from PostgreSQL itself. Both report that
restricting the schema shown to the model is not enough on its own, because models approximate an
answer from the columns that remain. We evaluate schemagate, an open-source library, which combines
three mechanisms: (i) the caller's permissions decide which tables and columns enter the prompt, so
restricted names are never shown; (ii) one line tells the model its view is partial, without naming
what is withheld; and (iii) every generated query is checked against the caller's permissions before
it runs. Scored by each benchmark's own protocol, it reaches an AC-F1 of 90.1 on BIRD in the RBAC
benchmark (best reported in that paper: 88.1, same model) and a refusal F1 of 0.984 on all 19,624
Spider-ACL records (best reported: 0.933). No query it answered read a table or column outside the
caller's grants, which we verify independently of the mechanism that enforces it. We also report the
negative results: on Spider the RBAC result ties the same model and trails the paper's best; adding
column descriptions made violations worse; and the first Spider-ACL pilot exposed a bug in which a
fail-closed query check silently refused valid SQL, which we argue is a general hazard of
fail-closed designs.
\end{abstract}

\section{Introduction}
A natural-language interface to a database inherits the database's access rules. A question about
salaries from a user without the payroll role should be refused, and refused before any data is read.
The common approaches are (a) to describe the policy to the model in the prompt and trust it to comply,
and (b) to show the model only the schema the caller may read. Fei et al.~\cite{fei2027rbac} find that
(a) leaves violations at several percent even for strong models, and that (b) fails in a different way:
the model, unaware anything is missing, answers anyway from the columns it can see. Klisura
et al.~\cite{klisura2026acl} report a similar trade-off between refusing correctly and answering
permitted questions.

This report evaluates a combination of (b) with two additions, on both benchmarks, using each
benchmark's own scoring code or scoring rule unchanged. Our contributions are measurements, not a new
model: no fine-tuning is involved.

\section{Method}
\paragraph{Identity-scoped schema selection.} The library reflects the database into a catalogue. Each
table, view and column carries the roles that may read it, either from a configuration or from the
database's own GRANTs. For a question and a caller, it selects the relevant objects from those the
caller may read and renders their DDL for the prompt. Restricted objects are absent rather than ranked
low.

\paragraph{Restriction notice.} When the caller cannot see some part of the catalogue, the rendered
schema begins with one line: \emph{``This caller may read only part of this database: some tables and
columns are not shown. Answer only if the tables and columns below directly hold what the question asks
for. If it would need data that is not shown, reply INSUFFICIENT; do not approximate it from other
columns.''} It names nothing that is withheld.

\paragraph{Scope check.} Before a generated query runs, the tables and columns it references are
extracted twice, by pattern readers and by a SQL parser (sqlglot) in the database's dialect, and the
query is refused if either finds a name the caller may not read, or if it cannot be parsed.
\texttt{SELECT *} over a table with withheld columns is refused.

\section{Benchmarks and setup}
\paragraph{RBAC Text-to-SQL~\cite{fei2027rbac}.} Spider~\cite{spider} and BIRD~\cite{bird} questions, each paired with a role and a
column-level policy. We ran the authors' evaluator unmodified with \texttt{--fair\_comparison
--num\_trials 5 --execute\_sql}, using claude-sonnet-4-5, the model behind the paper's best results.
Metrics: AC-F1, SafeEX, violation rate (any answer to a forbidden question), over-refusal, and the
rate of SQL that reads outside the policy.

\paragraph{Spider-ACL~\cite{klisura2026acl}.} 19,624 Spider question--role pairs; roles are PostgreSQL
GRANTs and the label is PERMIT or DENY as PostgreSQL returned when the reference query ran as that role
(11,529 DENY). We score refusal precision, recall and F1 with refusal as the positive class, following
the paper, using claude-sonnet-5-5 on all records.

\section{Results}
\begin{table}[h]\centering\small
\resizebox{\linewidth}{!}{\begin{tabular}{lccccc}\toprule
BIRD (10,175 rows) & AC-F1 & SafeEX & Violation & Over-refusal & Out-of-policy SQL\\\midrule
schemagate & \textbf{90.14}\,$\pm$\,0.32 & \textbf{68.38}\,$\pm$\,1.07 & \textbf{5.44}\,$\pm$\,0.19 & 1.98 & \textbf{0.00}\\
\quad without scope check & 84.69\,$\pm$\,0.69 & 69.14\,$\pm$\,0.85 & 11.25\,$\pm$\,0.44 & 1.23 & 3.78\\
Paper, same model, policy in prompt & 88.12\,$\pm$\,0.90 & 65.02\,$\pm$\,0.89 & 7.37\,$\pm$\,0.49 & 1.81 & --\\
Paper, best schema-only (GPT-5-mini) & 72.46 & -- & 26.72 & 0.33 & --\\\midrule
Spider (6,926 rows) & & & & & \\
schemagate & 93.02\,$\pm$\,0.37 & 78.87\,$\pm$\,1.05 & 5.05\,$\pm$\,0.49 & 2.59 & \textbf{0.00}\\
\quad without scope check & 90.68\,$\pm$\,0.31 & 79.56\,$\pm$\,1.05 & 8.67\,$\pm$\,0.57 & 1.91 & 2.01\\
Paper, same model & 93.03\,$\pm$\,0.48 & 80.99\,$\pm$\,0.86 & 4.55\,$\pm$\,0.22 & 3.01 & --\\
Paper, best (DeepSeek V3.2-Reasoning) & \textbf{95.94}\,$\pm$\,0.37 & 75.73\,$\pm$\,1.19 & \textbf{1.97}\,$\pm$\,0.25 & 2.36 & --\\\bottomrule
\end{tabular}}
\caption{RBAC Text-to-SQL benchmark, authors' evaluator, mean $\pm$ s.d.\ over five trials. On BIRD the
same-model result is also the best of the paper's thirteen models.}
\end{table}

On BIRD, the combination exceeds every configuration reported in~\cite{fei2027rbac} on AC-F1, SafeEX and
violations. On Spider it ties the same model and trails the paper's best model. Removing the scope check
raises the violation rate from 5.4\% to 11.3\% on BIRD and from 5.1\% to 8.7\% on Spider, and is the only one of our
configurations that produced out-of-policy SQL.

\begin{table}[h]\centering\small
\resizebox{\linewidth}{!}{\begin{tabular}{lcccc}\toprule
Spider-ACL & Refusal P & Refusal R & Refusal F1 & DENY answered\\\midrule
schemagate, all 19,624 records & 0.985 & 0.982 & \textbf{0.984} & 1.8\% (206/11,529)\\
schemagate, 600 records not used in development & 0.989 & 0.983 & 0.986 & 1.7\%\\
Paper, best (LoRA-tuned Mistral-7B) & 0.936 & 0.929 & 0.933 & --\\
Paper, best without fine-tuning (DeepSeek-R1 + verifier) & 0.898 & 0.868 & 0.883 & --\\\bottomrule
\end{tabular}}
\caption{Spider-ACL. Confusion counts for the full run: TP 11,323, FP 173, FN 206, TN 7,922.}
\end{table}

\paragraph{Independent audit.} Every one of the 8,128 queries answered on Spider-ACL was parsed with
sqlglot and its tables and columns compared with the record's GRANTs, separately from the scope check
that enforces the rule. None referenced an ungranted table or column. Two queries the parser could not
read (the model appended prose after the SQL) were checked by hand.

\section{Ablations and negative results}
\paragraph{Restriction notice.} On a 100-question BIRD pilot (one trial), with the scope check on,
violations fell from 16\% to 5\% and AC-F1 rose from 0.818 to 0.921 with the notice. Its wording was
chosen on that pilot, whose questions belong to the test set; re-scoring BIRD with those 100 questions
removed gives AC-F1 $90.08\pm0.25$, unchanged. The full runs were not repeated without the notice.

\paragraph{Column descriptions hurt.} Loading BIRD's column descriptions as comments raised violations
on the same pilot (16\% to 26\% without the notice): richer notes on the visible columns helped the
model construct substitutes for questions it should have refused.

\paragraph{A fail-closed bug.} The first Spider-ACL pilot (200 records) scored far lower than expected.
32 of its 33 wrong refusals had one cause: the pattern reader matched \texttt{table [AS] alias}
case-sensitively, so \texttt{FROM students AS T1} was read as two tables, and the unknown ``T1'' was
refused. Because the check fails closed, the bug caused no leak and no error, only refusals that looked
like caution. Our own tests had not caught it; a benchmark with externally produced labels did, within
two days. We think this generalises: a fail-closed guard converts its own defects into plausible
refusals, so its false-refusal rate must be measured, not only its leak rate.

\section{Limitations}
The models differ from most of those in the original papers: the RBAC comparison is same-model, but the
Spider-ACL systems of~\cite{klisura2026acl} used 2024--25 models, and part of that gap is the model.
On Spider-ACL we score only the answer-or-refuse decision, not execution accuracy. BIRD-ACL and the
RBAC benchmark's LiveSQLBench subset were not run. The BIRD databases used are the 2024-06 dev release;
\cite{fei2027rbac} used a 2025-11 release. Neither result has yet been reproduced by the benchmarks'
authors; both have been sent to them.

\section{Reproducibility}
Code, harnesses (\texttt{benchmarks/rbac\_text2sql.py}, \texttt{benchmarks/acl\_text2sql.py}), the
evaluator reports and every Spider-ACL decision with its SQL are public at
\url{https://github.com/ashishsinha1602/schemagate} (release 1.6.1, Apache-2.0).

\begin{thebibliography}{9}
\bibitem{fei2027rbac} Y.~Fei, Y.~Jiang, Y.~Yang, X.~Xiao. Benchmarking Text-to-SQL under Role-Based
Access Control. \emph{Proc.\ ACM Manag.\ Data} (SIGMOD 2027). arXiv:2607.22115.
\bibitem{klisura2026acl} \DJ.~Klisura, J.~Khoury, A.~Kundu, R.~Krishnan, A.~Rios. Role-Conditioned
Refusals: Evaluating Access Control Reasoning in Large Language Models. \emph{Findings of EACL 2026},
pp.~6018--6034.
\bibitem{spider} T.~Yu et al. Spider: A Large-Scale Human-Labeled Dataset for Complex and Cross-Domain
Semantic Parsing and Text-to-SQL Task. EMNLP 2018.
\bibitem{bird} J.~Li et al. Can LLM Already Serve as a Database Interface? A Big Bench for Large-Scale
Database Grounded Text-to-SQLs. NeurIPS 2023.
\end{thebibliography}
\end{document}
